\section{Upper bound on $R_r(d,n)$}\label{Sec:Rdn}
To simplify notation we will identify the vertex set of $K_N$ with $[N]=\{1,\dots,N\}$ and the vertex set of $K_N\Osq \cdots \Osq K_N$ with $[N]^d$.  Suppose that, for each $i\in[d]$ we have a graph $G_i$ and a subgraph $H_i\subseteq G_i$. For $a\in V(G_d)$,  we write $H_1\Osq\dots\Osq H_{d-1}\Osq a$ for the copy of $H_1\Osq\dots\Osq H_{d-1}\subseteq G_1\Osq\dots\Osq G_{d}$ with vertex set $V(H_1)\times\dots\times V(H_{d-1})\times \{a\}$ (thus all vertices have $a$ as their $d$th coordinate; we will usually omit the braces around $\{a\}$ for simplicity). We define subgraphs such as $a_1\Osq \cdots\Osq a_{d-1}\Osq H_{d}$ analogously.
We will sometimes refer to induced subgraphs by their vertex sets: thus, for an edge-coloured graph $G$ we  say that $S\subset V(G)$ {\em contains a monochromatic $K_k$} if there is some set $T\subset S$ such that the induced subgraph $G[T]$ is a monochromatic copy of $K_k$.
